\section{Metodologia}
\label{sec:method}

\subsection{Tipo de pesquisa}

Pesquisa aplicada de natureza quantitativa, com método experimental sobre
dados secundários (\textit{corpus} público). O delineamento é
\textit{descritivo-comparativo}: um \textit{pipeline} de amostragem é
aplicado, três alocações alternativas são contrastadas, e cinco estimadores
do parâmetro populacional são comparados quanto a viés e variância.

\subsection{População, universo e estratos}

A população alvo é o conjunto $\mathcal{U}$ de \textit{tweets} em português
brasileiro do ToLD-Br~\cite{leite2020}, $N = 21{,}000$ exemplos. A
distribuição populacional sob a regra de maioria ($v \geq 2$) é apresentada
na Tabela~\ref{tab:population}. O desbalanceamento extremo
(\textit{racism} a $0{,}16\%$, \textit{xenophobia} a $0{,}20\%$) motiva o
tratamento estratificado.

\begin{table}[!t]
\centering
\caption{Distribuição populacional do ToLD-Br ($N=21{,}000$, votos $\geq 2$).}
\label{tab:population}
\begin{tabular}{lrr}
\toprule
Categoria      & Positivos & \% da população \\
\midrule
\textit{obscene}    & 2{,}403 & 11{,}44\% \\
\textit{insult}     & 1{,}869 &  8{,}90\% \\
\textit{homophobia} &   176   &  0{,}84\% \\
\textit{misogyny}   &   133   &  0{,}63\% \\
\textit{xenophobia} &    42   &  0{,}20\% \\
\textit{racism}     &    33   &  0{,}16\% \\
\bottomrule
\end{tabular}
\end{table}

\begin{figure}[!t]
\centering
\includegraphics[width=.85\linewidth]{toldbr_label_counts.png}
\caption{Positivos por categoria sob a regra de maioria ($v \geq 2$).
A escala destaca o desbalanceamento entre as categorias mais comuns
(\textit{obscene}, \textit{insult}) e as caudas raras
(\textit{racism}, \textit{xenophobia}).}
\label{fig:label-counts}
\end{figure}

\paragraph{Definição dos estratos.} Como o esquema é multi-rótulo, adota-se a
regra \textit{raro-vence}: cada \textit{tweet} é alocado ao estrato
correspondente ao rótulo majoritário positivo mais raro entre os seis ---
ou ao estrato \textit{clean} se nenhum rótulo for majoritariamente positivo.
A ordem de prioridade do mais raro ao mais comum é
$\text{racism} \to \text{xenophobia} \to \text{misogyny} \to \text{homophobia}
\to \text{insult} \to \text{obscene}$. Esta partição é mutuamente exclusiva e
preserva integralmente a contagem populacional dos estratos raros.

\subsection{Variável-alvo e variável auxiliar}

\paragraph{Variável-alvo.} Define-se
$y_i = \sum_{j=1}^{6} v_{ij}$ ($\texttt{total\_severity}$), inteiro em
$[0, 18]$, interpretável como \textit{intensidade total de preocupação dos
anotadores} sobre o \textit{tweet}. A média populacional verdadeira é
$\bar{Y} = 0{,}8591$.

\paragraph{Variável auxiliar.} Adota-se $x = \texttt{word\_count}$ (contagem
de palavras), barata de calcular sobre a população completa. A correlação de
Pearson populacional é $\rho_{xy} \approx 0{,}011$ --- pequena, dominada pelo
estrato \textit{clean}; as correlações \emph{intra-estrato} são mais
substanciais (até $0{,}27$ em \textit{racism}). Esta escolha explora
deliberadamente um auxiliar fraco para examinar o regime em que o ganho
teórico da regressão sobre a média estratificada deveria desaparecer
(Cochran~\cite{cochran1977}~§7.4); um auxiliar mais informativo
(\texttt{caps\_ratio}, \textit{embeddings}) é avaliado como trabalho futuro.

\subsection{Cálculo do tamanho amostral}

Aplicando as equações~(\ref{eq:n0}) e~(\ref{eq:n}) com $z=1{,}96$, $p=0{,}5$,
$e=0{,}05$:
\begin{equation*}
  n_0 \approx 384{,}16; \qquad n = \frac{384{,}16}{1 + 383{,}16/21000}
            \approx 377{,}3 \;\Rightarrow\; n = 378.
\end{equation*}

\subsection{Estratégias de alocação comparadas}
\label{sec:method:allocation}

Três alocações de $n=378$ entre os estratos são comparadas:

\textbf{Proporcional (Bolfarine \& Bussab).} $n_h = n \cdot N_h/N$. Resulta
em $n_h = 1$ para os estratos \textit{racism} e \textit{xenophobia}.

\textbf{Cochran-por-estrato com censo (Cochran §5.5).} Aplicação
independente da equação~(\ref{eq:n}) a cada $N_h$, com censo se
$n_h \geq N_h$. Total $\approx 1{,}299$ exemplos --- útil como referência
alternativa que recupera todos os estratos raros.

\textbf{Ótima de Neyman (Cochran §5.5; Bolfarine \& Bussab §5.5).}
Equação~(\ref{eq:neyman}). Esta é a alocação adotada como referência
principal para a estimação de $\bar{Y}$.

Em qualquer das alocações, a amostragem dentro de cada estrato é
\textit{aleatória simples sem reposição} (SRS-WOR).

\subsection{Estimação do parâmetro populacional}

A partir da amostra Neyman, $\bar{Y}$ é estimado pelos cinco métodos descritos
na Seção~\ref{sec:theory:estimators}:
\begin{enumerate}
  \item Média estratificada simples (eq.~\ref{eq:strat-mean}).
  \item Estimador tipo razão separado (eq.~\ref{eq:ratio-sep}).
  \item Estimador tipo razão combinado (eq.~\ref{eq:ratio-comb}).
  \item Estimador tipo regressão separado (eq.~\ref{eq:reg-sep}).
  \item Estimador tipo regressão combinado (eq.~\ref{eq:reg-comb}).
\end{enumerate}
Cada estimativa é acompanhada de erro-padrão e intervalo de confiança de
$95\%$ (Wald, normal padrão).

\subsection{Instrumentos e procedimento}

O \textit{pipeline} é integralmente implementado em Python 3.12, gerenciado
via \texttt{uv}, com as bibliotecas \texttt{pandas}, \texttt{numpy},
\texttt{scipy}, \texttt{scikit-learn}, \texttt{matplotlib} e
\texttt{datasets}. O código-fonte está organizado em três módulos:
\begin{itemize}
  \item \texttt{toxicity\_analysis.sampling}: fórmula de Cochran e três
        alocações (proporcional, Cochran-por-estrato, Neyman) com SRS-WOR
        dentro de cada estrato.
  \item \texttt{toxicity\_analysis.features}: \textit{features} léxicas
        determinísticas e diagnóstico de ruído por contagem de votos.
  \item \texttt{toxicity\_analysis.estimators}: estimadores tipo razão e
        regressão, separado e combinado, com variâncias de grande amostra
        sob FPC.
\end{itemize}
A reprodutibilidade é garantida por uma suíte de $74$ testes unitários
(\texttt{pytest}) que cobre identidades textbook (e.g., $n_0 \approx 384$ a
$95\%$ IC; estimadores recuperam $\bar{Y}$ exatamente sob censo total) e
inputs adversariais. Todos os experimentos usam semente
\texttt{random\_state=42}.
